\section{Grover operator}\label{sec:groverOperator}

\red{We derive quantum inference schemes based on the Grover operator, which we want to first introduce in full generality.}

\subsection{Definition}

Let us assume we have two orthogonal states
\begin{itemize}
    \item $\qstateofat{\goodstatesymbol}{\shortcatvariables}$
    \item $\qstateofat{\badstatesymbol}{\shortcatvariables}$
\end{itemize}
spanning a real subspace $\subspaceof{}=\spanof{\qstateofat{\goodstatesymbol}{\shortcatvariables},\qstateofat{\goodstatesymbol}{\shortcatvariables}}$ and let us choose another state
\begin{align*}
    \qstateofat{0}{\shortcatvariables}
    \in \subspaceof{}\, .
\end{align*}
Then there is $\rotanglesymbol\in[0,4\pi)$ with
\begin{align*}
    \qstateofat{0}{\shortcatvariables}
    = \sinof{\frac{\rotanglesymbol}{2}}\qstateofat{\goodstatesymbol}{\shortcatvariables}
    + \cosof{\frac{\rotanglesymbol}{2}}\qstateofat{\badstatesymbol}{\shortcatvariables} \, .
\end{align*}
Consider the reflection about the states $\qstateof{\badstatesymbol}$ and $\qstateof{0}$ by
\begin{align*}
    \reflectionofat{\badstatesymbol}{\inshortcatvariables,\outshortcatvariables}
    &= 2 \cdot \qstateofat{\badstatesymbol}{\inshortcatvariables} \otimes \qstateofat{\badstatesymbol}{\outshortcatvariables} - \identityat{\inshortcatvariables,\outshortcatvariables} \\
    \reflectionofat{0}{\inshortcatvariables,\outshortcatvariables}
    &= 2 \cdot \qstateofat{0}{\inshortcatvariables} \otimes \qstateofat{0}{\outshortcatvariables} - \identityat{\inshortcatvariables,\outshortcatvariables} \, ,
\end{align*}
and their concatenation
\begin{align*}
    \groverat{\inshortcatvariables,\outshortcatvariables}
    = \contractionof{
        \reflectionofat{\badstatesymbol}{\inshortcatvariables,\ccatvariableof{[\catorder]}{\midsymbol}},
        \reflectionofat{0}{\ccatvariableof{[\catorder]}{\midsymbol},\outshortcatvariables}
    }{\inshortcatvariables,\outshortcatvariables}
\end{align*}
which is called the \emph{Grover operator}.

\subsection{Action on projected quantum states}

The Grover operator leaves any quantum state $\qstate$ invariant, which is orthogonal to $\qstateof{\goodstatesymbol}$ and $\qstateof{\badstatesymbol}$.
The projection of the quantum state onto $\subspaceof{}$ gets rotated by the Grover operator (see \figref{fig:rotationSketch}).
To be more precise, we denote the projection by $(\mathcal{P}\qstate)[\selvariable]$ which has the coordinates
\begin{align*}
    \left(\mathcal{P}_{\subspaceof{}}\qstate\right)\left[\selvariable=0\right]
    &= \contraction{\left(\qstatewith\right)^{\dagger},\qstateofat{\goodstatesymbol}{\shortcatvariables}}\\
    \left(\mathcal{P}_{\subspaceof{}}\qstate\right)\left[\selvariable=1\right]
    &= \contraction{\left(\qstatewith\right)^{\dagger},\qstateofat{\goodstatesymbol}{\shortcatvariables}} \, .
\end{align*}
We define a matrix representation of the Grover operator as
\begin{align*}
    \groverofat{\mathrm{red}}{\cselvariable{\outsymbol},\cselvariable{\insymbol}}
    = \coloredmatrixof{
        \cosof{\rotanglesymbol} & -\sinof{\rotanglesymbol} \\
        \sinof{\rotanglesymbol} & \cosof{\rotanglesymbol}
    }{\cselvariable{\outsymbol},\cselvariable{\insymbol}}
\end{align*}
and it holds for any quantum state $\qstatewith$ that
\begin{align*}
    \left(\mathcal{P}_{\subspaceof{}}\groversymbol\qstate\right)\left[\cselvariable{\outsymbol}\right]
    = \contractionof{\groverofat{\mathrm{red}}{\cselvariable{\outsymbol},\cselvariable{\insymbol}},
        \left(\mathcal{P}_{\subspaceof{}}\qstate\right)\left[\cselvariable{\insymbol}\right]
    }{\outsymbol} \, .
\end{align*}
The Grover operator therefore has the eigenvalues:
\begin{itemize}
    \item $\expof{\pm i \rotanglesymbol}$ with the eigen spaces spanned by $\qstateof{\goodstatesymbol}\pm i \qstateof{\badstatesymbol}$
    \item $1$ with the eigen space by the orthogonal complement of $\subspaceof{}$
\end{itemize}
Application the Grover operator $\repnum\in\nn$ times on $\qstateof{0}$ results in the state
\begin{align*}
    \qstateofat{\repnum}{\shortcatvariables}
    = \sinof{\left(\frac{1}{2}+\repnum\right)\rotanglesymbol} \qstateofat{\goodstatesymbol}{\shortcatvariables}
    + \cosof{\left(\frac{1}{2}+\repnum\right)\rotanglesymbol} \qstateofat{\badstatesymbol}{\shortcatvariables} \, .
\end{align*}

\begin{figure}[t]
    \begin{center}
        \input{partII/5-quantum-inference/tikz/grover-operator/rotation_sketch.tikz}
    \end{center}
    \caption{Grover operator as a rotation in $\spanof{\qstateof{\goodstatesymbol},\qstateof{\badstatesymbol}}$.}
    \label{fig:rotationSketch}
\end{figure}

\subsection{Circuit Realization}

Let us now investigate how the reflection operations can be realized by circuits.
The reflection $\reflectionof{}$ about the ground state $\onehotmapofat{\zerotuple{[\catorder]}}{\shortcatvariables}$ is realized by Pauli-X and controlled Pauli-Z.
Having a unitary $\unitarysymbol$ preparing $\qstateof{0}$ by action on the ground state, we realize the reflection
\begin{align*}
    \reflectionofat{0}{\inshortcatvariables,\outshortcatvariables}
    = \contractionof{
        \unitaryofat{\dagger}{\inshortcatvariables,\ccatvariableof{[\catorder]}{\midsymbol,0}},
        \reflectionofat{}{\ccatvariableof{[\catorder]}{\midsymbol,0},\ccatvariableof{[\catorder]}{\midsymbol,1}},
        \unitaryat{\ccatvariableof{[\catorder]}{\midsymbol,1},\outshortcatvariables}
    }{\inshortcatvariables,\outshortcatvariables} \, .
\end{align*}

We will investigate different effective schemes to realize $\reflectionof{\badstatesymbol}$ later.



